\documentclass[11pt,a4paper]{article}

% ============================================================
%  QIMMA -- COURS : Limites et continuité
%  2ème Bac Sciences Mathématiques (A et B)
%  Niveau renforcé SM : définitions formelles + démonstrations
%  Basé sur templates/qimma-template.tex (charte Qimma)
% ============================================================

\usepackage[french]{babel}
\usepackage[utf8]{inputenc}
\usepackage[T1]{fontenc}
\usepackage{lmodern}
\usepackage[a4paper,margin=2cm,top=2.3cm,bottom=2.6cm]{geometry}
\usepackage{amsmath,amssymb}
\usepackage{xcolor}
\usepackage{graphicx}
\usepackage{tikz}
\usetikzlibrary{shadows,calc}
\usepackage[most]{tcolorbox}
\usepackage{titlesec}
\usepackage{fancyhdr}
\usepackage{enumitem}
\usepackage{pifont}
\usepackage{array}
\usepackage{colortbl}
\usepackage{hyperref}

% ------------------- CHARTE QIMMA -------------------
\definecolor{qteal}{HTML}{0d9488}
\definecolor{qtealdark}{HTML}{134e4a}
\definecolor{qamber}{HTML}{fbbf24}
\definecolor{qambertext}{HTML}{92400e}
\definecolor{ink}{HTML}{1f2937}
\definecolor{inkgray}{HTML}{6b7280}
\definecolor{tealpale}{HTML}{e6f5f3}
\colorlet{colDef}{qteal}
\definecolor{colProp}{HTML}{0f766e}
\definecolor{colEx}{HTML}{b45309}
\definecolor{colRem}{HTML}{9d174d}
\color{ink}

\graphicspath{{../../../../assets/}}
\newcommand{\logofile}{../../../../assets/qimma-logo.pdf}

% ------------------- METADONNEES -------------------
\newcommand{\matiere}{Mathématiques}
\newcommand{\niveau}{2\up{ème} Bac -- Sciences Mathématiques}
\newcommand{\chapitretitre}{Limites et continuité}

% ------------------- EN-TETE / PIED DE PAGE -------------------
\pagestyle{fancy}
\fancyhf{}
\renewcommand{\headrulewidth}{0pt}
\renewcommand{\footrulewidth}{0.5pt}
\fancyfoot[L]{\raisebox{-0.15cm}{\includegraphics[height=0.35cm]{\logofile}}~\small\sffamily\color{inkgray}Qimma}
\fancyfoot[C]{\small\sffamily\color{inkgray}\matiere\ -- \niveau}
\fancyfoot[R]{\small\sffamily\color{qtealdark}\thepage}

% ------------------- TITRES DE SECTION -------------------
\titleformat{\section}
  {\normalfont\sffamily\Large\bfseries\color{qtealdark}}
  {\tikz[baseline=-3.2pt]{\node[circle,fill=qteal,text=white,inner sep=0pt,minimum size=8mm]{\sffamily\bfseries\thesection};}}
  {10pt}{}
  [\vspace{-2mm}{\color{qamber}\titlerule[1.6pt]}\vspace{2mm}]
\titlespacing*{\section}{0pt}{16pt}{10pt}

\titleformat{\subsection}
  {\normalfont\sffamily\large\bfseries\color{qtealdark}}
  {\color{qteal}\ding{212}}{6pt}{}
\titlespacing*{\subsection}{0pt}{12pt}{6pt}

% ------------------- BOITES PEDAGOGIQUES -------------------
\newtcolorbox{fiche}[3][]{
  enhanced, breakable,
  colback=white, colframe=#2,
  boxrule=1pt, arc=2mm,
  top=7mm, left=4mm, right=4mm, bottom=3mm,
  drop shadow={black!25, shadow xshift=1mm, shadow yshift=-1mm},
  attach boxed title to top left={xshift=4mm, yshift=-3mm, yshifttext=-1mm},
  boxed title style={colback=#2, arc=1.5mm, boxrule=0pt, left=2mm, right=2mm},
  coltitle=white, fonttitle=\sffamily\bfseries,
  title={#3}, #1
}
\newenvironment{definition}[1][Définition]
  {\begin{fiche}{colDef}{\ding{110}~~#1}}{\end{fiche}}
\newenvironment{propriete}[1][Propriété]
  {\begin{fiche}{colProp}{\ding{72}~~#1}}{\end{fiche}}
\newenvironment{exemple}[1][Exemple]
  {\begin{fiche}{colEx}{\ding{217}~~#1}}{\end{fiche}}
\newenvironment{remarque}[1][Remarque]
  {\begin{fiche}{colRem}{\ding{43}~~#1}}{\end{fiche}}

% Démonstration (hors boîte, style discret)
\newenvironment{demo}
  {\par\smallskip\noindent{\sffamily\bfseries\color{qtealdark}Démonstration.}\ \itshape}
  {\ \normalfont\hfill$\blacksquare$\par\medskip}

\newcommand{\E}{\mathrm{E}}

\hypersetup{colorlinks=true, linkcolor=qtealdark, urlcolor=qtealdark}

% ============================================================
\begin{document}

% ------------------- BANNIERE DE TITRE -------------------
\begin{tcolorbox}[
  enhanced, boxrule=0pt, arc=4mm,
  interior style={left color=qtealdark, right color=qteal},
  top=8mm, bottom=8mm, left=10mm, right=32mm,
  drop shadow={black!35, shadow xshift=1.5mm, shadow yshift=-1.5mm},
  before skip=0pt, after skip=9mm,
  overlay={
    \node[anchor=north west] at ([xshift=6mm,yshift=-6mm]frame.north west)
      {\includegraphics[height=1.25cm]{\logofile}};
    \node[anchor=north east, fill=qamber, text=qambertext, rounded corners=2pt,
          inner sep=4pt, font=\sffamily\bfseries\small] at ([xshift=-6mm,yshift=-6mm]frame.north east) {COURS};
  }
]
\hspace{1.6cm}\centering
{\sffamily\bfseries\color{white}\fontsize{23}{27}\selectfont \chapitretitre}\\[3mm]
{\sffamily\color{white!90}\large \matiere\ \textbullet\ \niveau}
\end{tcolorbox}

% ------------------- OBJECTIFS -------------------
\begin{tcolorbox}[
  enhanced, colback=tealpale, colframe=qteal, boxrule=0.8pt, arc=2mm,
  left=5mm, right=5mm, top=3mm, bottom=3mm,
  title={\sffamily\bfseries\color{qtealdark}\ding{51}~~Objectifs du chapitre},
  coltitle=qtealdark, colbacktitle=tealpale, boxrule=0.8pt,
  attach title to upper={\par\vspace{1mm}}
]
\begin{itemize}[leftmargin=6mm, label=\ding{212}, itemsep=1mm]
  \item Manipuler la définition formelle (en $\varepsilon$) de la limite et l'utiliser dans des démonstrations simples.
  \item Calculer des limites en levant les formes indéterminées (terme dominant, conjugué, limites trigonométriques, encadrement avec la partie entière).
  \item Étudier la continuité en un point, sur un intervalle, et construire un prolongement par continuité.
  \item Exploiter le lien entre suites et continuité, notamment pour les suites récurrentes $u_{n+1} = f(u_n)$.
  \item Démontrer l'existence et l'unicité de solutions d'équations par le théorème des valeurs intermédiaires, et encadrer une solution par dichotomie avec une précision donnée.
  \item Appliquer le théorème de la bijection, étudier une fonction réciproque, et maîtriser la fonction racine $n$-ième et les puissances rationnelles.
\end{itemize}
\end{tcolorbox}

\vspace{4mm}

% ------------------- SECTION 1 -------------------
\section{La notion de limite : définitions formelles}

\subsection{Limite finie en un point}

\begin{definition}[Limite finie en un point (définition en $\varepsilon$)]
Soit $f$ une fonction définie sur un intervalle ouvert contenant $a$, sauf peut-être en $a$, et soit $\ell \in \mathbb{R}$. On dit que $f$ \textbf{admet $\ell$ pour limite en $a$}, et on écrit $\lim\limits_{x \to a} f(x) = \ell$, lorsque :
\[
  (\forall \varepsilon > 0)\ (\exists \alpha > 0)\ (\forall x \in D_f)\quad
  0 < |x - a| < \alpha \implies |f(x) - \ell| < \varepsilon.
\]
Autrement dit : on peut rendre $f(x)$ aussi proche de $\ell$ que l'on veut, à condition de prendre $x$ suffisamment proche de $a$.
\end{definition}

\begin{propriete}[Unicité de la limite]
Si $f$ admet une limite en $a$, cette limite est \textbf{unique}.
\end{propriete}

\begin{demo}
Supposons que $f$ admette deux limites $\ell \neq \ell'$ en $a$, et posons $\varepsilon = \frac{|\ell - \ell'|}{2} > 0$. Il existe $\alpha_1 > 0$ tel que $0 < |x - a| < \alpha_1 \Rightarrow |f(x) - \ell| < \varepsilon$, et $\alpha_2 > 0$ tel que $0 < |x - a| < \alpha_2 \Rightarrow |f(x) - \ell'| < \varepsilon$. Pour $0 < |x - a| < \min(\alpha_1, \alpha_2)$, l'inégalité triangulaire donne
\[
  |\ell - \ell'| \leq |\ell - f(x)| + |f(x) - \ell'| < 2\varepsilon = |\ell - \ell'|,
\]
ce qui est absurde. Donc $\ell = \ell'$.
\end{demo}

\begin{exemple}[Utiliser la définition en $\varepsilon$]
Montrons avec la définition que $\lim\limits_{x \to 2} (3x - 1) = 5$.

\smallskip
Soit $\varepsilon > 0$. On a $|(3x - 1) - 5| = |3x - 6| = 3|x - 2|$. Ainsi, en choisissant $\alpha = \dfrac{\varepsilon}{3}$ :
\[
  0 < |x - 2| < \alpha \implies |(3x - 1) - 5| = 3|x - 2| < 3\alpha = \varepsilon.
\]
La définition est vérifiée : $\lim\limits_{x \to 2}(3x - 1) = 5$.
\end{exemple}

\subsection{Limites infinies, limites en l'infini, limites à droite et à gauche}

\begin{definition}[Limites infinies et limites en l'infini]
Avec les mêmes conventions :
\begin{itemize}[leftmargin=6mm, itemsep=0.5mm]
  \item $\lim\limits_{x \to a} f(x) = +\infty$ lorsque : $(\forall A > 0)(\exists \alpha > 0)$ : $0 < |x - a| < \alpha \implies f(x) > A$.
  \item $\lim\limits_{x \to +\infty} f(x) = \ell$ lorsque : $(\forall \varepsilon > 0)(\exists B > 0)$ : $x > B \implies |f(x) - \ell| < \varepsilon$.
  \item $\lim\limits_{x \to +\infty} f(x) = +\infty$ lorsque : $(\forall A > 0)(\exists B > 0)$ : $x > B \implies f(x) > A$.
\end{itemize}
Les définitions en $-\infty$ et pour $-\infty$ s'écrivent de façon analogue.
\end{definition}

\begin{definition}[Limites à droite et à gauche]
$\lim\limits_{x \to a^+} f(x) = \ell$ (limite \textbf{à droite}) lorsque la condition de la définition est exigée seulement pour $a < x < a + \alpha$ ; de même à gauche avec $a - \alpha < x < a$.
\[
  \lim_{x \to a} f(x) = \ell
  \iff
  \lim_{x \to a^-} f(x) = \lim_{x \to a^+} f(x) = \ell.
\]
\end{definition}

\begin{exemple}[Limites latérales distinctes]
Soit $f(x) = \dfrac{|x - 1|}{x - 1}$ pour $x \neq 1$. Pour $x > 1$, $f(x) = 1$ ; pour $x < 1$, $f(x) = -1$. Donc
\[
  \lim_{x \to 1^+} f(x) = 1 \neq -1 = \lim_{x \to 1^-} f(x) :
\]
$f$ n'admet pas de limite en $1$.
\end{exemple}

\begin{remarque}[La fonction partie entière]
La fonction \textbf{partie entière} $\E$ associe à $x$ l'unique entier $\E(x)$ tel que $\E(x) \leq x < \E(x) + 1$. En tout point entier $n$, elle vérifie $\lim\limits_{x \to n^-} \E(x) = n - 1$ et $\lim\limits_{x \to n^+} \E(x) = n$ : elle n'a pas de limite en $n$ (et n'y est pas continue). L'encadrement fondamental $x - 1 < \E(x) \leq x$ est l'outil clé pour toutes les limites faisant intervenir $\E$.
\end{remarque}

% ------------------- SECTION 2 -------------------
\section{Calcul des limites : opérations, ordre, composition}

\subsection{Opérations et composition}

\begin{propriete}[Opérations sur les limites]
Si $\lim f = \ell$ et $\lim g = \ell'$ (limites finies, au même endroit), alors :
\[
  \lim (f + g) = \ell + \ell', \qquad
  \lim (fg) = \ell\,\ell', \qquad
  \lim \frac{f}{g} = \frac{\ell}{\ell'} \ (\ell' \neq 0),
\]
\[
  \lim |f| = |\ell|, \qquad
  \lim \sqrt{f} = \sqrt{\ell} \ (\ell \geq 0).
\]
Les règles s'étendent aux limites infinies avec les conventions usuelles ($\ell + \infty = \infty$, etc.), \textbf{sauf} pour les quatre \textbf{formes indéterminées} :
\[
  \text{« } \infty - \infty \text{ »}, \qquad
  \text{« } 0 \times \infty \text{ »}, \qquad
  \text{« } \tfrac{\infty}{\infty} \text{ »}, \qquad
  \text{« } \tfrac{0}{0} \text{ »}.
\]
\end{propriete}

\begin{propriete}[Limite d'une fonction composée]
Si $\lim\limits_{x \to a} f(x) = b$ et $\lim\limits_{t \to b} g(t) = c$ (où $a$, $b$, $c$ sont finis ou infinis), alors
\[
  \lim_{x \to a} (g \circ f)(x) = c.
\]
\end{propriete}

\begin{exemple}[Composée]
Calculons $\displaystyle\lim_{x \to +\infty} \sqrt{\frac{4x^2 + 1}{x^2 + 3}}$.

\smallskip
D'abord la limite interne (forme $\frac{\infty}{\infty}$, on factorise par $x^2$) :
\[
  \frac{4x^2 + 1}{x^2 + 3} = \frac{4 + \frac{1}{x^2}}{1 + \frac{3}{x^2}}
  \xrightarrow[x \to +\infty]{} 4.
\]
Puis, par composition avec $t \mapsto \sqrt{t}$ continue en $4$ : la limite cherchée vaut $\sqrt{4} = 2$.
\end{exemple}

\subsection{Limites et ordre}

\begin{propriete}[Théorèmes de comparaison et d'encadrement]
Au voisinage considéré (en $\pm\infty$ ou en un point $a$) :
\begin{enumerate}[leftmargin=7mm, itemsep=0.5mm]
  \item si $f \geq g$ et $g \to +\infty$, alors $f \to +\infty$ ; si $f \leq g$ et $g \to -\infty$, alors $f \to -\infty$ ;
  \item \textbf{(encadrement, « théorème des gendarmes »)} si $g \leq f \leq h$ et si $g$ et $h$ tendent vers la \emph{même} limite finie $\ell$, alors $f \to \ell$ ;
  \item si $|f - \ell| \leq g$ et $g \to 0$, alors $f \to \ell$ ;
  \item \textbf{(passage à la limite dans une inégalité)} si $f \leq g$ et si les deux limites existent, alors $\lim f \leq \lim g$.
\end{enumerate}
\end{propriete}

\begin{demo}
Démontrons le point 3 en un point $a$. Soit $\varepsilon > 0$. Comme $g \to 0$, il existe $\alpha > 0$ tel que $0 < |x - a| < \alpha \Rightarrow |g(x)| < \varepsilon$. Pour ces $x$ :
\[
  |f(x) - \ell| \leq g(x) \leq |g(x)| < \varepsilon,
\]
ce qui est exactement la définition de $\lim\limits_{x \to a} f(x) = \ell$. Le point 2 s'en déduit en remarquant que $|f - \ell| \leq \max(|g - \ell|, |h - \ell|) \to 0$.
\end{demo}

\begin{exemple}[Limite avec la partie entière -- niveau SM]
Calculons $\displaystyle\lim_{x \to 0^+} x\,\E\!\left(\frac{1}{x}\right)$.

\smallskip
Pour $x > 0$, l'encadrement $\frac{1}{x} - 1 < \E\!\left(\frac{1}{x}\right) \leq \frac{1}{x}$, multiplié par $x > 0$, donne
\[
  1 - x < x\,\E\!\left(\frac{1}{x}\right) \leq 1.
\]
Quand $x \to 0^+$, $1 - x \to 1$ : par le théorème des gendarmes,
$\displaystyle\lim_{x \to 0^+} x\,\E\!\left(\frac{1}{x}\right) = 1$.
\end{exemple}

\subsection{Limites trigonométriques}

\begin{propriete}[Limites trigonométriques usuelles]
\[
  \lim_{x \to 0} \frac{\sin x}{x} = 1, \qquad
  \lim_{x \to 0} \frac{\tan x}{x} = 1, \qquad
  \lim_{x \to 0} \frac{1 - \cos x}{x^2} = \frac{1}{2}.
\]
\end{propriete}

\begin{demo}
On admet l'encadrement géométrique classique, valable pour $0 < x < \frac{\pi}{2}$ (comparaison d'aires dans le cercle trigonométrique) : $\sin x < x < \tan x$. En divisant par $\sin x > 0$ : $1 < \frac{x}{\sin x} < \frac{1}{\cos x}$, d'où $\cos x < \frac{\sin x}{x} < 1$. Comme $\cos x \to 1$ en $0$, les gendarmes donnent $\frac{\sin x}{x} \to 1$ (la fonction étant paire, la limite vaut aussi à gauche).

Ensuite $\frac{\tan x}{x} = \frac{\sin x}{x} \cdot \frac{1}{\cos x} \to 1 \times 1 = 1$.

Enfin, avec $1 - \cos x = 2\sin^2\frac{x}{2}$ :
\[
  \frac{1 - \cos x}{x^2} = \frac{2\sin^2\frac{x}{2}}{x^2}
  = \frac{1}{2}\left(\frac{\sin\frac{x}{2}}{\frac{x}{2}}\right)^{\!2}
  \xrightarrow[x \to 0]{} \frac{1}{2} \times 1^2 = \frac{1}{2}.
\]
\end{demo}

\begin{exemple}[Deux calculs types]
\textbf{a)} $\displaystyle\lim_{x \to 0} \frac{\sin 5x}{\sin 2x}$ : on écrit
\[
  \frac{\sin 5x}{\sin 2x}
  = \frac{\sin 5x}{5x} \cdot \frac{2x}{\sin 2x} \cdot \frac{5}{2}
  \xrightarrow[x \to 0]{} 1 \times 1 \times \frac{5}{2} = \frac{5}{2}.
\]
\textbf{b)} $\displaystyle\lim_{x \to 0} \frac{1 - \cos x}{x \sin x}$ : on écrit
\[
  \frac{1 - \cos x}{x \sin x}
  = \frac{1 - \cos x}{x^2} \cdot \frac{x}{\sin x}
  \xrightarrow[x \to 0]{} \frac{1}{2} \times 1 = \frac{1}{2}.
\]
\end{exemple}

\begin{exemple}[Conjugué + racine]
Calculons $\displaystyle\lim_{x \to +\infty} \left(\sqrt{x^2 + x} - x\right)$ (forme « $\infty - \infty$ »).
\[
  \sqrt{x^2 + x} - x
  = \frac{(x^2 + x) - x^2}{\sqrt{x^2 + x} + x}
  = \frac{x}{x\left(\sqrt{1 + \frac{1}{x}} + 1\right)}
  = \frac{1}{\sqrt{1 + \frac{1}{x}} + 1}
  \xrightarrow[x \to +\infty]{} \frac{1}{2}.
\]
\emph{Attention en $-\infty$} : pour $x < 0$, $\sqrt{x^2} = |x| = -x$ ; toute factorisation sous une racine doit en tenir compte.
\end{exemple}

\subsection{Lever les formes indéterminées : toutes les méthodes}

\begin{propriete}[Polynômes et fractions rationnelles en l'infini]
En $+\infty$ et en $-\infty$, un polynôme a la limite de son \textbf{terme de plus haut degré}, et une fraction rationnelle la limite du \textbf{quotient des termes de plus haut degré} :
\[
  \lim_{x \to \pm\infty} \bigl(a_n x^n + \cdots + a_0\bigr) = \lim_{x \to \pm\infty} a_n x^n,
  \qquad
  \lim_{x \to \pm\infty} \frac{a_n x^n + \cdots}{b_m x^m + \cdots} = \lim_{x \to \pm\infty} \frac{a_n x^n}{b_m x^m}.
\]
(Résultat obtenu en factorisant par $x^n$ et $x^m$ ; à \textbf{redémontrer} par factorisation si l'énoncé l'exige.)
\end{propriete}

\begin{exemple}[Rationnelle en l'infini : les trois cas]
\textbf{a)} Degrés égaux : $\displaystyle\lim_{x \to +\infty} \frac{2x^2 - 5x + 1}{3x^2 + 4} = \lim_{x \to +\infty}\frac{2x^2}{3x^2} = \frac{2}{3}$.

\smallskip
\textbf{b)} Degré du haut plus grand : $\displaystyle\lim_{x \to -\infty} \frac{x^3 + 1}{x + 2} = \lim_{x \to -\infty}\frac{x^3}{x} = \lim_{x \to -\infty} x^2 = +\infty$.

\smallskip
\textbf{c)} Degré du bas plus grand : $\displaystyle\lim_{x \to +\infty} \frac{5x + 3}{x^2 + 1} = \lim_{x \to +\infty}\frac{5x}{x^2} = \lim_{x \to +\infty}\frac{5}{x} = 0$.
\end{exemple}

\begin{exemple}[Forme $0/0$ rationnelle : factoriser la racine commune]
Calculons $\displaystyle\lim_{x \to 3} \frac{x^2 - 9}{x^2 - 5x + 6}$ (numérateur et dénominateur s'annulent en $3$).
\[
  \frac{x^2 - 9}{x^2 - 5x + 6}
  = \frac{(x - 3)(x + 3)}{(x - 3)(x - 2)}
  = \frac{x + 3}{x - 2}
  \xrightarrow[x \to 3]{} \frac{6}{1} = 6.
\]
\emph{Règle} : quand $a$ annule les deux membres d'un quotient de polynômes, $(x - a)$ se factorise en haut et en bas.
\end{exemple}

\begin{exemple}[En $-\infty$ avec racine : le piège du signe]
Calculons $\displaystyle\lim_{x \to -\infty} \left(\sqrt{x^2 + 1} + x\right)$ (forme « $\infty - \infty$ » car $\sqrt{x^2+1} \to +\infty$ et $x \to -\infty$).

\smallskip
Conjugué :
\[
  \sqrt{x^2 + 1} + x = \frac{(x^2 + 1) - x^2}{\sqrt{x^2 + 1} - x} = \frac{1}{\sqrt{x^2 + 1} - x}.
\]
Quand $x \to -\infty$ : $-x \to +\infty$ et $\sqrt{x^2+1} \to +\infty$, donc le dénominateur tend vers $+\infty$ et la \textbf{limite vaut $0$}.

\smallskip
\emph{Variante par factorisation} : pour $x < 0$, $\sqrt{x^2 + 1} = |x|\sqrt{1 + \frac{1}{x^2}} = -x\sqrt{1 + \frac{1}{x^2}}$, d'où $\sqrt{x^2+1} + x = -x\left(\sqrt{1 + \frac{1}{x^2}} - 1\right)$ : produit d'un facteur $\to +\infty$ par un facteur $\to 0$, à traiter alors par conjugué. La factorisation avec le \textbf{signe correct} de $|x|$ est le point le plus surveillé par les correcteurs.
\end{exemple}

\begin{exemple}[Changement de variable]
Calculons $\displaystyle\lim_{x \to +\infty} x \sin\frac{1}{x}$ (forme « $\infty \times 0$ »).

\smallskip
Posons $t = \frac{1}{x}$ : quand $x \to +\infty$, $t \to 0^+$, et
\[
  x \sin\frac{1}{x} = \frac{\sin t}{t} \xrightarrow[t \to 0^+]{} 1.
\]
\end{exemple}

\begin{remarque}[Tableau des réflexes selon la forme rencontrée]
\begin{center}
\renewcommand{\arraystretch}{1.4}
\sffamily\small
\begin{tabular}{>{\raggedright\arraybackslash}p{6.2cm} >{\raggedright\arraybackslash}p{8.6cm}}
\rowcolor{qtealdark} \color{white}\bfseries Situation & \color{white}\bfseries Méthode \\
\rowcolor{tealpale} Polynôme ou rationnelle en $\pm\infty$ & Termes de plus haut degré (factorisation). \\
« $0/0$ » avec des polynômes & Factoriser la racine commune $(x - a)$. \\
\rowcolor{tealpale} « $\infty - \infty$ » ou « $0/0$ » avec $\sqrt{\phantom{x}}$ & Expression conjuguée (attention au signe de $|x|$ en $-\infty$). \\
« $0/0$ » trigonométrique en $0$ & Se ramener à $\frac{\sin X}{X}$, $\frac{\tan X}{X}$, $\frac{1 - \cos X}{X^2}$. \\
\rowcolor{tealpale} « $0 \times \infty$ » & Réécrire en quotient ou changement de variable $t = \frac{1}{x}$. \\
Présence de $\E(x)$ & Encadrement $x - 1 < \E(x) \leq x$ puis gendarmes. \\
\rowcolor{tealpale} Présence de $\cos$, $\sin$ bornés & Encadrer entre $-1$ et $1$ puis gendarmes ou point 3. \\
Racine $n$-ième en « $0/0$ » & Changement de variable $t = \sqrt[n]{x}$ ou conjugué cubique. \\
\end{tabular}
\end{center}
\end{remarque}

% ------------------- SECTION 3 -------------------
\section{Continuité en un point et prolongement par continuité}

\begin{definition}[Continuité en un point]
Soit $f$ définie sur un intervalle ouvert contenant $a$. On dit que $f$ est \textbf{continue en $a$} lorsque $\lim\limits_{x \to a} f(x) = f(a)$, c'est-à-dire :
\[
  (\forall \varepsilon > 0)\ (\exists \alpha > 0)\ (\forall x \in D_f)\quad
  |x - a| < \alpha \implies |f(x) - f(a)| < \varepsilon.
\]
$f$ est \textbf{continue à droite} en $a$ si $\lim\limits_{x \to a^+} f(x) = f(a)$, \textbf{continue à gauche} si $\lim\limits_{x \to a^-} f(x) = f(a)$ ; elle est continue en $a$ si et seulement si elle l'est à droite et à gauche.
\end{definition}

\begin{propriete}[Opérations]
Si $f$ et $g$ sont continues en $a$, alors $f + g$, $\lambda f$, $fg$, $|f|$ sont continues en $a$ ; $\frac{f}{g}$ l'est si $g(a) \neq 0$ ; $\sqrt{f}$ l'est si $f(a) > 0$ (ou $f \geq 0$ au voisinage). Si $f$ est continue en $a$ et $g$ continue en $f(a)$, alors $g \circ f$ est continue en $a$.

\smallskip
Conséquences : les polynômes sont continus sur $\mathbb{R}$, les fonctions rationnelles sur leur domaine, $\sin$, $\cos$ sur $\mathbb{R}$, $\tan$ sur son domaine, $x \mapsto \sqrt{x}$ sur $[0\,;+\infty[$.
\end{propriete}

\begin{definition}[Prolongement par continuité]
Soit $f$ définie sur un intervalle ouvert contenant $a$, \textbf{sauf en $a$}, telle que $\lim\limits_{x \to a} f(x) = \ell \in \mathbb{R}$. La fonction $\tilde{f}$ définie par
\[
  \tilde{f}(x) =
  \begin{cases}
    f(x) & \text{si } x \neq a,\\[1mm]
    \ell & \text{si } x = a,
  \end{cases}
\]
est continue en $a$ : on l'appelle le \textbf{prolongement par continuité} de $f$ en $a$.
\end{definition}

\begin{exemple}[Construire un prolongement par continuité]
Soit $f(x) = \dfrac{\sqrt{x + 1} - 1}{x}$, définie sur $]-1\,;0[\ \cup\ ]0\,;+\infty[$. $f$ est-elle prolongeable par continuité en $0$ ?

\smallskip
Multiplions par le conjugué :
\[
  f(x) = \frac{(x + 1) - 1}{x\left(\sqrt{x+1} + 1\right)} = \frac{1}{\sqrt{x + 1} + 1}
  \xrightarrow[x \to 0]{} \frac{1}{2}.
\]
La limite existe et est finie : $f$ est prolongeable par continuité en $0$, avec $\tilde{f}(0) = \frac{1}{2}$.
\end{exemple}

\begin{exemple}[Continuité d'une fonction définie par morceaux]
Soit $f(x) = \dfrac{\sin x}{x}$ si $x \neq 0$ et $f(0) = 1$. Comme $\lim\limits_{x \to 0} \frac{\sin x}{x} = 1 = f(0)$, $f$ est continue en $0$ ; ailleurs elle est continue comme quotient de fonctions continues à dénominateur non nul. Conclusion : $f$ est continue sur $\mathbb{R}$.
\end{exemple}

\begin{exemple}[Déterminer des paramètres pour rendre $f$ continue]
Soit $f$ définie sur $\mathbb{R}$ par :
\[
  f(x) =
  \begin{cases}
    x + a & \text{si } x < 0,\\
    x^2 + 1 & \text{si } 0 \leq x \leq 1,\\
    bx + 2 & \text{si } x > 1.
  \end{cases}
\]
Déterminons $a$ et $b$ pour que $f$ soit continue sur $\mathbb{R}$.

\smallskip
Sur chacun des trois intervalles ouverts, $f$ est polynomiale donc continue : seuls les points de raccordement $0$ et $1$ sont à étudier.

\smallskip
\textbf{En $0$} : $f(0) = 1$ ; $\lim\limits_{x \to 0^-} f(x) = a$ ; $\lim\limits_{x \to 0^+} f(x) = 1$. Continuité en $0$ $\iff a = 1$.

\smallskip
\textbf{En $1$} : $f(1) = 2$ ; $\lim\limits_{x \to 1^-} f(x) = 2$ ; $\lim\limits_{x \to 1^+} f(x) = b + 2$. Continuité en $1$ $\iff b + 2 = 2 \iff b = 0$.

\smallskip
\textbf{Conclusion} : $a = 1$ et $b = 0$.
\end{exemple}

% ------------------- SECTION 4 -------------------
\section{Suites et continuité}

\begin{propriete}[Image d'une suite convergente]
Soit $(u_n)$ une suite telle que $u_n \to a$, et $f$ une fonction \textbf{continue en $a$} (avec $u_n \in D_f$). Alors
\[
  \lim_{n \to +\infty} f(u_n) = f(a).
\]
\end{propriete}

\begin{propriete}[Limite d'une suite récurrente]
Soit $(u_n)$ définie par $u_{n+1} = f(u_n)$. Si $(u_n)$ \textbf{converge} vers $\ell$ et si $f$ est \textbf{continue en $\ell$}, alors $\ell$ est un \textbf{point fixe} de $f$ :
\[
  f(\ell) = \ell.
\]
\end{propriete}

\begin{demo}
Puisque $u_n \to \ell$ et $f$ est continue en $\ell$, on a $f(u_n) \to f(\ell)$. Or $f(u_n) = u_{n+1}$, et la suite $(u_{n+1})$ converge vers la même limite $\ell$ que $(u_n)$. Par unicité de la limite, $f(\ell) = \ell$.
\end{demo}

\begin{exemple}[Schéma type bac SM]
Soit $u_0 = 1$ et $u_{n+1} = \sqrt{u_n + 2}$. On suppose démontré (par récurrence) que $0 \leq u_n \leq 2$ et que $(u_n)$ est croissante. Justifions la convergence et déterminons la limite.

\smallskip
$(u_n)$ est croissante et majorée par $2$ : elle converge vers un réel $\ell \in [0\,;2]$ (théorème de la convergence monotone). La fonction $f(t) = \sqrt{t + 2}$ est continue sur $[0\,;2]$, donc $\ell$ vérifie $\ell = \sqrt{\ell + 2}$, soit $\ell^2 - \ell - 2 = 0$, c'est-à-dire $(\ell - 2)(\ell + 1) = 0$. Comme $\ell \geq 0$ : $\boxed{\ell = 2}$.
\end{exemple}

\begin{remarque}[Erreur à éviter]
L'équation $f(\ell) = \ell$ ne donne la limite \emph{que si l'on a déjà prouvé la convergence}. Résoudre le point fixe sans justifier la convergence (monotonie + borne, ou $|f'| \leq k < 1$) ne rapporte aucun point.
\end{remarque}

% ------------------- SECTION 5 -------------------
\section{Continuité sur un intervalle et théorème des valeurs intermédiaires}

\begin{definition}
$f$ est \textbf{continue sur} $]a\,;b[$ si elle est continue en tout point de $]a\,;b[$ ; \textbf{continue sur} $[a\,;b]$ si elle est de plus continue à droite en $a$ et à gauche en $b$.
\end{definition}

\begin{propriete}[{Image d'un intervalle, image d'un segment}]
L'image d'un \textbf{intervalle} par une fonction continue est un \textbf{intervalle}. L'image d'un \textbf{segment} $[a\,;b]$ est un \textbf{segment} $[m\,;M]$, où $m$ et $M$ sont le minimum et le maximum de $f$ sur $[a\,;b]$ (une fonction continue sur un segment est bornée et atteint ses bornes).
\end{propriete}

\begin{remarque}[Lecture directe pour $f$ strictement monotone]
$f$ croissante : $f([a\,;b]) = [f(a)\,;f(b)]$ ; $f\bigl(\,]a\,;b]\,\bigr) = \bigl]\lim_{a^+} f\,;f(b)\bigr]$, etc. -- on remplace chaque borne par la valeur ou la limite correspondante, en respectant le caractère ouvert/fermé. $f$ décroissante : les bornes s'échangent.
\end{remarque}

\begin{propriete}[Théorème des valeurs intermédiaires (TVI)]
Soit $f$ \textbf{continue} sur $[a\,;b]$. Pour tout $k$ compris entre $f(a)$ et $f(b)$, il existe au moins un $c \in [a\,;b]$ tel que $f(c) = k$.

\smallskip
\textbf{Corollaire 1} : si $f$ est continue sur $[a\,;b]$ et $f(a)\,f(b) < 0$, l'équation $f(x) = 0$ admet au moins une solution dans $]a\,;b[$.

\smallskip
\textbf{Corollaire 2 (unicité)} : si de plus $f$ est \textbf{strictement monotone} sur $[a\,;b]$, cette solution est unique.
\end{propriete}

\begin{demo}
Idée de la preuve du corollaire 1 par \textbf{dichotomie}. Supposons $f(a) < 0 < f(b)$. On construit deux suites $(a_n)$ et $(b_n)$ : on part de $a_0 = a$, $b_0 = b$ ; à chaque étape, on considère le milieu $m_n = \frac{a_n + b_n}{2}$ : si $f(m_n) \leq 0$ on pose $a_{n+1} = m_n$, $b_{n+1} = b_n$, sinon $a_{n+1} = a_n$, $b_{n+1} = m_n$. Par construction $f(a_n) \leq 0 \leq f(b_n)$ et $b_n - a_n = \frac{b - a}{2^n} \to 0$ : les suites sont adjacentes et convergent vers un même réel $c$. Par continuité de $f$ et passage à la limite dans les inégalités, $f(c) \leq 0$ et $f(c) \geq 0$, donc $f(c) = 0$.
\end{demo}

\begin{exemple}[TVI + dichotomie avec précision imposée]
Montrons que l'équation $x^3 + x - 3 = 0$ admet une unique solution $\alpha$ dans $]1\,;2[$, puis donnons un encadrement de $\alpha$ d'amplitude $0{,}25$.

\smallskip
Posons $f(x) = x^3 + x - 3$, continue sur $[1\,;2]$ (polynôme), avec $f'(x) = 3x^2 + 1 > 0$ donc $f$ strictement croissante ; $f(1) = -1 < 0$ et $f(2) = 7 > 0$. Par le TVI et la stricte monotonie, $\alpha$ existe et est unique dans $]1\,;2[$.

\smallskip
\textbf{Dichotomie} : $f(1{,}5) = 3{,}375 + 1{,}5 - 3 = 1{,}875 > 0 \Rightarrow \alpha \in\, ]1\,;1{,}5[$. Puis $f(1{,}25) \approx 1{,}953 + 1{,}25 - 3 = 0{,}203 > 0 \Rightarrow \alpha \in\, ]1\,;1{,}25[$.

\smallskip
\textbf{Conclusion} : $1 < \alpha < 1{,}25$, encadrement d'amplitude $0{,}25$ obtenu en $2$ étapes (l'amplitude après $n$ étapes est $\frac{b - a}{2^n}$ ; pour une précision $10^{-2}$ il faudrait $\frac{1}{2^n} \leq 10^{-2}$, soit $n = 7$ étapes).
\end{exemple}

\begin{remarque}[Rédaction attendue à l'examen]
Pour appliquer le TVI, cite \emph{toujours} : \textbf{continuité} sur l'intervalle, \textbf{valeurs aux bornes} (ou changement de signe), et pour l'unicité la \textbf{stricte monotonie}. Précise ensuite l'intervalle ouvert ou fermé où vit la solution.
\end{remarque}

\begin{propriete}[Variantes du TVI à connaître absolument]
\begin{enumerate}[leftmargin=7mm, itemsep=0.5mm]
  \item \textbf{Équation $f(x) = k$} : appliquer le TVI à $g(x) = f(x) - k$, ou vérifier directement que $k$ est compris entre $f(a)$ et $f(b)$.
  \item \textbf{Équation $g(x) = h(x)$} : poser $f = g - h$ et chercher un changement de signe de $f$.
  \item \textbf{Intervalle non borné} : si $f$ est continue et strictement croissante sur $[a\,;+\infty[$ avec $f(a) < 0$ et $\lim\limits_{x \to +\infty} f(x) = +\infty$, alors $f(x) = 0$ admet une unique solution dans $]a\,;+\infty[$ (on remplace la valeur à la borne par la \textbf{limite}).
  \item \textbf{Polynôme de degré impair} : tout polynôme de degré impair admet \textbf{au moins une racine réelle}.
\end{enumerate}
\end{propriete}

\begin{demo}
Point 4 : soit $P$ de degré impair, de coefficient dominant $a_n > 0$ (sinon considérer $-P$). Alors $\lim\limits_{x \to -\infty} P(x) = -\infty$ et $\lim\limits_{x \to +\infty} P(x) = +\infty$ : il existe donc $a$ avec $P(a) < 0$ et $b$ avec $P(b) > 0$. $P$ étant continu sur $[a\,;b]$, le TVI fournit une racine dans $]a\,;b[$.
\end{demo}

\begin{exemple}[Équation $g(x) = h(x)$ sur un intervalle non borné]
Montrons que l'équation $x^3 = 3 - x$ admet une unique solution réelle $\alpha$, puis que $\alpha \in\, ]1\,;\frac{3}{2}[$.

\smallskip
Posons $f(x) = x^3 + x - 3$ sur $\mathbb{R}$ : $f$ est continue (polynôme) et $f'(x) = 3x^2 + 1 > 0$, donc $f$ est strictement croissante sur $\mathbb{R}$. De plus $\lim\limits_{x \to -\infty} f(x) = -\infty$ et $\lim\limits_{x \to +\infty} f(x) = +\infty$ : $f$ réalise une bijection de $\mathbb{R}$ sur $\mathbb{R}$, donc $f(x) = 0$ admet une \textbf{unique} solution $\alpha$.

\smallskip
Localisation : $f(1) = -1 < 0$ et $f\!\left(\frac{3}{2}\right) = \frac{27}{8} + \frac{3}{2} - 3 = \frac{15}{8} > 0$, donc $\alpha \in\, \bigl]1\,;\frac{3}{2}\bigr[$.
\end{exemple}

% ------------------- SECTION 6 -------------------
\section{Fonction réciproque d'une fonction continue strictement monotone}

\begin{propriete}[Théorème de la bijection]
Soit $f$ \textbf{continue et strictement monotone} sur un intervalle $I$. Alors :
\begin{enumerate}[leftmargin=7mm, itemsep=0.5mm]
  \item $f$ réalise une \textbf{bijection} de $I$ sur l'intervalle $J = f(I)$ ;
  \item la \textbf{réciproque} $f^{-1} : J \to I$, définie par $\bigl(y = f(x) \iff x = f^{-1}(y)\bigr)$, est continue sur $J$ et strictement monotone, de \textbf{même sens de variation} que $f$ ;
  \item dans un repère \textbf{orthonormé}, les courbes $\mathcal{C}_f$ et $\mathcal{C}_{f^{-1}}$ sont symétriques par rapport à la première bissectrice $y = x$.
\end{enumerate}
De plus, pour tout $x \in I$ : $f^{-1}(f(x)) = x$, et pour tout $y \in J$ : $f(f^{-1}(y)) = y$.
\end{propriete}

\begin{demo}
Démontrons le sens de variation (cas $f$ strictement croissante). Soient $y_1 < y_2$ dans $J$, et $x_1 = f^{-1}(y_1)$, $x_2 = f^{-1}(y_2)$. Si l'on avait $x_1 \geq x_2$, la croissance stricte de $f$ donnerait $y_1 = f(x_1) \geq f(x_2) = y_2$, contradiction. Donc $x_1 < x_2$ : $f^{-1}$ est strictement croissante. (La continuité de $f^{-1}$ est admise.)
\end{demo}

\begin{exemple}[Étude complète d'une réciproque]
Soit $f(x) = x^2 + 2x$ sur $I = [-1\,;+\infty[$.

\smallskip
\textbf{Bijection} : $f$ est dérivable et $f'(x) = 2x + 2 = 2(x + 1) \geq 0$, nulle seulement en $-1$ : $f$ est continue et strictement croissante sur $I$. Or $f(-1) = -1$ et $\lim\limits_{x \to +\infty} f(x) = +\infty$, donc $f$ réalise une bijection de $[-1\,;+\infty[$ sur $J = [-1\,;+\infty[$.

\smallskip
\textbf{Expression de $f^{-1}$} : pour $y \geq -1$, on résout $y = x^2 + 2x$ avec $x \geq -1$. On écrit $y = (x + 1)^2 - 1$, d'où $(x + 1)^2 = y + 1$ et, comme $x + 1 \geq 0$ : $x = -1 + \sqrt{y + 1}$.
\[
  f^{-1} : [-1\,;+\infty[\ \to\ [-1\,;+\infty[, \qquad f^{-1}(y) = \sqrt{y + 1} - 1.
\]
\textbf{Vérification} : $f\bigl(f^{-1}(y)\bigr) = (\sqrt{y+1} - 1)^2 + 2(\sqrt{y+1} - 1) = (y + 1) - 2\sqrt{y+1} + 1 + 2\sqrt{y+1} - 2 = y$. $\checkmark$
\end{exemple}

% ------------------- SECTION 7 -------------------
\section{Racine $n$-ième et puissances rationnelles}

\begin{definition}[Fonction racine $n$-ième]
Soit $n \in \mathbb{N}^*$. La fonction $x \mapsto x^n$ est continue et strictement croissante sur $[0\,;+\infty[$, donc bijective de $[0\,;+\infty[$ sur $[0\,;+\infty[$. Sa réciproque est la \textbf{fonction racine $n$-ième} :
\[
  y = \sqrt[n]{x} \iff y^n = x \qquad (x \geq 0,\ y \geq 0).
\]
Elle est continue et strictement croissante sur $[0\,;+\infty[$, avec $\lim\limits_{x \to +\infty} \sqrt[n]{x} = +\infty$.
\end{definition}

\begin{propriete}[Équation $x^n = a$]
\begin{itemize}[leftmargin=6mm, itemsep=0.5mm]
  \item $n$ \textbf{impair} : pour tout $a \in \mathbb{R}$, l'équation $x^n = a$ admet une \textbf{unique} solution réelle, notée $\sqrt[n]{a}$ (avec $\sqrt[n]{a} = -\sqrt[n]{|a|}$ si $a < 0$).
  \item $n$ \textbf{pair} : si $a > 0$, deux solutions $\pm\sqrt[n]{a}$ ; si $a = 0$, la seule solution est $0$ ; si $a < 0$, \textbf{aucune} solution réelle.
\end{itemize}
\end{propriete}

\begin{propriete}[Règles de calcul]
Pour tous $a, b \geq 0$ et $n, p \in \mathbb{N}^*$, $m \in \mathbb{N}$ :
\[
  \sqrt[n]{ab} = \sqrt[n]{a}\,\sqrt[n]{b}, \qquad
  \sqrt[n]{\frac{a}{b}} = \frac{\sqrt[n]{a}}{\sqrt[n]{b}}\ (b \neq 0), \qquad
  \sqrt[n]{\sqrt[p]{a}} = \sqrt[np]{a},
\]
\[
  \bigl(\sqrt[n]{a}\bigr)^m = \sqrt[n]{a^m}, \qquad
  \sqrt[np]{a^{mp}} = \sqrt[n]{a^m}.
\]
\end{propriete}

\begin{demo}
Montrons la première : posons $u = \sqrt[n]{a}$ et $v = \sqrt[n]{b}$, donc $u, v \geq 0$, $u^n = a$, $v^n = b$. Alors $(uv)^n = u^n v^n = ab$ et $uv \geq 0$ : par \emph{unicité} de la racine $n$-ième de $ab$, $uv = \sqrt[n]{ab}$. Les autres règles se démontrent sur le même schéma (élever à la bonne puissance et invoquer l'unicité).
\end{demo}

\begin{definition}[Puissances rationnelles]
Pour $a > 0$ et $r = \frac{p}{q} \in \mathbb{Q}$ ($p \in \mathbb{Z}$, $q \in \mathbb{N}^*$) :
\[
  a^{r} = a^{p/q} = \sqrt[q]{a^p}.
\]
Les règles usuelles des puissances restent valables : $a^r a^{r'} = a^{r + r'}$, $(a^r)^{r'} = a^{rr'}$, $(ab)^r = a^r b^r$.
\end{definition}

\begin{exemple}[Calculs et équation]
\textbf{a)} Simplifions $A = \sqrt[3]{16} \times \sqrt[3]{4}$ : $A = \sqrt[3]{64} = 4$.

\smallskip
\textbf{b)} Simplifions $B = \dfrac{\sqrt[4]{81} \times \sqrt{3}}{\sqrt[4]{9}}$ : on écrit tout en base $3$ : $\sqrt[4]{81} = 3$, $\sqrt{3} = 3^{1/2}$, $\sqrt[4]{9} = 3^{2/4} = 3^{1/2}$. Donc $B = \dfrac{3 \times 3^{1/2}}{3^{1/2}} = 3$.

\smallskip
\textbf{c)} Résolvons $x^5 = -32$ : $n = 5$ impair, solution unique $x = -\sqrt[5]{32} = -2$.

\smallskip
\textbf{d)} Résolvons $\sqrt[3]{x - 1} = 2$ : en élevant au cube, $x - 1 = 8$, soit $x = 9$ (pas de condition de signe à vérifier, le cube étant bijectif sur $\mathbb{R}$).
\end{exemple}

\begin{exemple}[Limite avec racine $n$-ième]
Calculons $\displaystyle\lim_{x \to 1} \frac{\sqrt[3]{x} - 1}{x - 1}$ (forme « $0/0$ »).

\smallskip
Posons $t = \sqrt[3]{x}$, donc $x = t^3$ et $t \to 1$ quand $x \to 1$ :
\[
  \frac{\sqrt[3]{x} - 1}{x - 1}
  = \frac{t - 1}{t^3 - 1}
  = \frac{t - 1}{(t - 1)(t^2 + t + 1)}
  = \frac{1}{t^2 + t + 1}
  \xrightarrow[t \to 1]{} \frac{1}{3}.
\]
\end{exemple}

\begin{propriete}[{Ordre, équations et limites avec les racines $n$-ièmes}]
\begin{itemize}[leftmargin=6mm, itemsep=0.5mm]
  \item \textbf{Stricte croissance} : pour $a, b \geq 0$ : $a < b \iff \sqrt[n]{a} < \sqrt[n]{b}$. On résout donc les équations et inéquations en \textbf{élevant à la puissance $n$} (après contrôle du domaine).
  \item \textbf{Continuité et limites} : la fonction $\sqrt[n]{\phantom{x}}$ est continue sur $[0\,;+\infty[$ ; si $f \geq 0$ et $\lim f = \ell$ (fini ou $+\infty$), alors $\lim \sqrt[n]{f} = \sqrt[n]{\ell}$ (avec $\sqrt[n]{+\infty} = +\infty$).
  \item \textbf{Identité du conjugué cubique} : $a^3 - b^3 = (a - b)(a^2 + ab + b^2)$, qui permet de lever les « $0/0$ » contenant des racines cubiques.
\end{itemize}
\end{propriete}

\begin{exemple}[Inéquations avec racines $n$-ièmes]
\textbf{a)} Résolvons $\sqrt[3]{2x - 1} \leq 3$. La racine cubique est définie sur $\mathbb{R}$ et strictement croissante : l'inéquation équivaut à $2x - 1 \leq 27$, soit $x \leq 14$. $S = \,]-\infty\,;14]$.

\smallskip
\textbf{b)} Résolvons $\sqrt[4]{x} > 2$. \textbf{Domaine} : $x \geq 0$. Par stricte croissance, l'inéquation équivaut à $x > 2^4 = 16$. $S = \,]16\,;+\infty[$.
\end{exemple}

\begin{exemple}[Conjugué cubique]
Calculons $\displaystyle\lim_{x \to 0} \frac{\sqrt[3]{1 + x} - 1}{x}$ (forme « $0/0$ »).

\smallskip
Avec $a = \sqrt[3]{1+x}$ et $b = 1$, on multiplie par $a^2 + ab + b^2$ :
\[
  \frac{\sqrt[3]{1+x} - 1}{x}
  = \frac{(1 + x) - 1}{x\left(\sqrt[3]{(1+x)^2} + \sqrt[3]{1+x} + 1\right)}
  = \frac{1}{\sqrt[3]{(1+x)^2} + \sqrt[3]{1+x} + 1}
  \xrightarrow[x \to 0]{} \frac{1}{3}.
\]
\end{exemple}

% ------------------- SECTION 8 -------------------
\section{Exercices corrigés -- niveau examen national SM}

\begin{exemple}[Exercice 1 -- limite trigonométrique complète]
Calculer $\displaystyle\lim_{x \to 0} \frac{\sqrt{1 + \sin x} - \sqrt{1 - \sin x}}{x}$.

\smallskip
\textbf{Corrigé.} Conjugué au numérateur :
\[
  \frac{(1 + \sin x) - (1 - \sin x)}{x\left(\sqrt{1 + \sin x} + \sqrt{1 - \sin x}\right)}
  = \frac{2\sin x}{x\left(\sqrt{1 + \sin x} + \sqrt{1 - \sin x}\right)}
  = \frac{\sin x}{x} \cdot \frac{2}{\sqrt{1 + \sin x} + \sqrt{1 - \sin x}}.
\]
Quand $x \to 0$ : $\frac{\sin x}{x} \to 1$ et le second facteur tend vers $\frac{2}{1 + 1} = 1$. \textbf{Limite : $1$.}
\end{exemple}

\begin{exemple}[Exercice 2 -- continuité avec paramètre]
Soit $f(x) = \dfrac{x^2 - 4}{x - 2}$ si $x \neq 2$ et $f(2) = m$. Pour quelle valeur du paramètre $m$ la fonction $f$ est-elle continue en $2$ ?

\smallskip
\textbf{Corrigé.} Pour $x \neq 2$ : $f(x) = \dfrac{(x-2)(x+2)}{x - 2} = x + 2$, donc $\lim\limits_{x \to 2} f(x) = 4$. $f$ est continue en $2$ si et seulement si $m = 4$.
\end{exemple}

\begin{exemple}[Exercice 3 -- TVI sur une équation non polynomiale]
Montrer que l'équation $\cos x = x$ admet une unique solution dans $\left[0\,;\frac{\pi}{2}\right]$.

\smallskip
\textbf{Corrigé.} Posons $g(x) = \cos x - x$, continue sur $\left[0\,;\frac{\pi}{2}\right]$ (différence de fonctions continues). $g$ est dérivable et $g'(x) = -\sin x - 1 < 0$ sur $\left]0\,;\frac{\pi}{2}\right[$ : $g$ est strictement décroissante. De plus $g(0) = 1 > 0$ et $g\!\left(\frac{\pi}{2}\right) = -\frac{\pi}{2} < 0$. Par le TVI et la stricte monotonie, l'équation $g(x) = 0$, c'est-à-dire $\cos x = x$, admet une unique solution dans $\left]0\,;\frac{\pi}{2}\right[$.
\end{exemple}

\begin{exemple}[Exercice 4 -- bijection et image d'intervalle]
Soit $f(x) = \dfrac{2x}{x + 1}$ sur $I = \,]-1\,;+\infty[$. Montrer que $f$ réalise une bijection de $I$ sur un intervalle $J$ à déterminer, et expliciter $f^{-1}$.

\smallskip
\textbf{Corrigé.} $f$ est dérivable sur $I$ et
\[
  f'(x) = \frac{2(x + 1) - 2x}{(x + 1)^2} = \frac{2}{(x + 1)^2} > 0 :
\]
$f$ est continue et strictement croissante sur $I$. Aux bornes : $\lim\limits_{x \to (-1)^+} f(x) = -\infty$ (numérateur $\to -2$, dénominateur $\to 0^+$) et $\lim\limits_{x \to +\infty} f(x) = 2$. Donc $f$ est une bijection de $I$ sur $J = \,]-\infty\,;2[$.

\smallskip
Pour $y \in J$, résolvons $y = \frac{2x}{x+1}$ : $y(x + 1) = 2x \iff x(y - 2) = -y \iff x = \frac{y}{2 - y}$ (licite car $y \neq 2$).
\[
  f^{-1} : \,]-\infty\,;2[\ \to\ ]-1\,;+\infty[, \qquad f^{-1}(y) = \frac{y}{2 - y}.
\]
\end{exemple}

\begin{exemple}[Exercice 5 -- limite en $-\infty$ avec racine (piège du signe)]
Calculer $\displaystyle\lim_{x \to -\infty} \left(\sqrt{x^2 - x + 1} + x\right)$.

\smallskip
\textbf{Corrigé.} Forme « $\infty - \infty$ ». Conjugué :
\[
  \sqrt{x^2 - x + 1} + x = \frac{(x^2 - x + 1) - x^2}{\sqrt{x^2 - x + 1} - x}
  = \frac{-x + 1}{\sqrt{x^2 - x + 1} - x}.
\]
Pour $x < 0$, on factorise par $-x > 0$ : au numérateur $-x + 1 = -x\left(1 - \frac{1}{x}\right)$ ; au dénominateur $\sqrt{x^2 - x + 1} = -x\sqrt{1 - \frac{1}{x} + \frac{1}{x^2}}$, donc
\[
  \frac{-x\left(1 - \frac{1}{x}\right)}{-x\left(\sqrt{1 - \frac{1}{x} + \frac{1}{x^2}} + 1\right)}
  = \frac{1 - \frac{1}{x}}{\sqrt{1 - \frac{1}{x} + \frac{1}{x^2}} + 1}
  \xrightarrow[x \to -\infty]{} \frac{1}{1 + 1} = \frac{1}{2}.
\]
\textbf{Limite : $\dfrac{1}{2}$.}
\end{exemple}

\begin{exemple}[Exercice 6 -- limite avec partie entière]
Calculer $\displaystyle\lim_{x \to +\infty} \frac{\E(x)}{x}$.

\smallskip
\textbf{Corrigé.} Pour $x > 0$, l'encadrement $x - 1 < \E(x) \leq x$ donne, en divisant par $x > 0$ :
\[
  1 - \frac{1}{x} < \frac{\E(x)}{x} \leq 1.
\]
Comme $1 - \frac{1}{x} \to 1$, le théorème des gendarmes donne $\displaystyle\lim_{x \to +\infty} \frac{\E(x)}{x} = 1$.
\end{exemple}

\begin{exemple}[Exercice 7 -- bijection avec racine cubique]
Soit $f(x) = \sqrt[3]{x - 1} + 2$. Montrer que $f$ est une bijection de $\mathbb{R}$ sur $\mathbb{R}$ et expliciter $f^{-1}$.

\smallskip
\textbf{Corrigé.} $f$ est la composée de $x \mapsto x - 1$ (continue, strictement croissante sur $\mathbb{R}$), de la racine cubique (continue, strictement croissante sur $\mathbb{R}$) et d'une translation : $f$ est continue et strictement croissante sur $\mathbb{R}$. Ses limites en $\mp\infty$ valent $\mp\infty$, donc $f(\mathbb{R}) = \mathbb{R}$ : $f$ est une bijection de $\mathbb{R}$ sur $\mathbb{R}$.

\smallskip
Pour $y \in \mathbb{R}$ : $y = \sqrt[3]{x - 1} + 2 \iff \sqrt[3]{x - 1} = y - 2 \iff x - 1 = (y - 2)^3$, d'où
\[
  f^{-1}(y) = (y - 2)^3 + 1.
\]
\end{exemple}

\begin{exemple}[Exercice 8 -- TVI et dichotomie avec précision imposée]
Montrer que l'équation $\sin x = 1 - x$ admet une unique solution $\alpha$ dans $[0\,;1]$, puis donner un encadrement de $\alpha$ d'amplitude $0{,}25$.

\smallskip
\textbf{Corrigé.} Posons $g(x) = \sin x + x - 1$, continue sur $[0\,;1]$. $g$ est dérivable et $g'(x) = \cos x + 1 > 0$ sur $]0\,;1[$ (car $\cos x > 0$ sur cet intervalle) : $g$ est strictement croissante. $g(0) = -1 < 0$ et $g(1) = \sin 1 \approx 0{,}84 > 0$. Par le TVI et la stricte monotonie, $\alpha$ existe et est unique dans $]0\,;1[$.

\smallskip
\textbf{Dichotomie} : $g(0{,}5) = \sin(0{,}5) - 0{,}5 \approx 0{,}479 - 0{,}5 < 0 \Rightarrow \alpha \in\, ]0{,}5\,;1[$. Puis $g(0{,}75) \approx 0{,}682 - 0{,}25 > 0 \Rightarrow \alpha \in\, ]0{,}5\,;0{,}75[$ : encadrement d'amplitude $0{,}25$.
\end{exemple}

% ------------------- SECTION 9 -------------------
\section{Méthodes et pièges : la boîte à outils de l'examen}

\subsection{Ce qu'on te demande, ce que tu fais}

\begin{center}
\renewcommand{\arraystretch}{1.45}
\sffamily\small
\begin{tabular}{>{\raggedright\arraybackslash}p{7.2cm} >{\raggedright\arraybackslash}p{8cm}}
\rowcolor{qtealdark} \color{white}\bfseries Ce qu'on te demande & \color{white}\bfseries Ton réflexe \\
\rowcolor{tealpale} Limite d'une rationnelle en $\pm\infty$ & Quotient des termes de plus haut degré. \\
Limite « $0/0$ » avec racines carrées & Conjugué ($a^2 - b^2$) ; cubiques : $a^3 - b^3$. \\
\rowcolor{tealpale} Limite trigonométrique en $0$ & Se ramener à $\frac{\sin X}{X} \to 1$. \\
Limite avec $\E(x)$, $\cos$, $\sin$ & Encadrement puis gendarmes. \\
\rowcolor{tealpale} Continuité en un point de raccordement & Limites à droite et à gauche $= f(a)$. \\
Prolongement par continuité en $a$ & Calculer $\lim_{x \to a} f$ ; si finie, poser $\tilde{f}(a) = \ell$. \\
\rowcolor{tealpale} Solution unique de $f(x) = k$ & TVI (continuité + valeurs/limites aux bornes) + stricte monotonie. \\
Encadrer $\alpha$ à la précision $p$ & Dichotomie : amplitude $\frac{b - a}{2^n} \leq p$. \\
\rowcolor{tealpale} Montrer que $f$ est une bijection de $I$ sur $J$ & Continuité + stricte monotonie + calcul de $J = f(I)$ par valeurs/limites. \\
Expliciter $f^{-1}$ & Résoudre $y = f(x)$ d'inconnue $x$, en gardant la contrainte $x \in I$. \\
\rowcolor{tealpale} Limite d'une suite $u_{n+1} = f(u_n)$ & 1) Prouver la convergence. 2) Résoudre $f(\ell) = \ell$. \\
Équation/inéquation avec $\sqrt[n]{\phantom{x}}$ & Domaine d'abord, puis élever à la puissance $n$ (stricte croissance). \\
\end{tabular}
\end{center}

\subsection{Les pièges qui coûtent des points}

\begin{remarque}[Erreurs relevées chaque année par les correcteurs]
\begin{itemize}[leftmargin=6mm, itemsep=0.5mm]
  \item En $-\infty$ : écrire $\sqrt{x^2} = x$ au lieu de $\sqrt{x^2} = -x$ (pour $x < 0$).
  \item Conclure par le point fixe $f(\ell) = \ell$ \textbf{sans avoir prouvé la convergence} de la suite.
  \item Appliquer le TVI sans citer la \textbf{continuité}, ou conclure à l'unicité sans \textbf{stricte monotonie}.
  \item Prolonger par continuité alors que la limite est \textbf{infinie} ou n'existe pas : impossible.
  \item Oublier le \textbf{domaine} avant d'élever une inéquation à une puissance paire ($\sqrt[4]{x} > 2$ exige $x \geq 0$).
  \item Écrire $\sqrt[n]{a^n} = a$ pour $a < 0$ et $n$ pair : c'est $|a|$.
  \item Pour un quotient tendant vers « $\frac{\ell}{0}$ » : oublier d'étudier le \textbf{signe du dénominateur} à droite et à gauche (limites $+\infty$ ou $-\infty$ différentes).
  \item Confondre $f^{-1}(x)$ et $\frac{1}{f(x)}$.
  \item Traiter « $\infty - \infty$ » ou « $\frac{\infty}{\infty}$ » comme si cela valait $0$ ou $1$ : ce sont des \textbf{formes indéterminées}.
  \item Étudier la continuité de $\E$ en un entier sans séparer limite à gauche ($n - 1$) et à droite ($n$).
\end{itemize}
\end{remarque}

% ------------------- RESUME -------------------
\vspace{4mm}
\begin{tcolorbox}[
  enhanced, boxrule=0pt, arc=2mm,
  interior style={left color=qtealdark, right color=qteal},
  left=6mm, right=6mm, top=4mm, bottom=4mm,
  fontupper=\color{white}\sffamily,
  title={\sffamily\bfseries\color{white}\ding{72}~~À retenir},
  colbacktitle=qtealdark, coltitle=white, boxrule=0pt
]
\begin{itemize}[leftmargin=6mm, label={\color{qamber}\ding{212}}, itemsep=1mm]
  \item Définition formelle : $\lim_{x\to a} f = \ell \iff (\forall\varepsilon>0)(\exists\alpha>0):\ 0<|x-a|<\alpha \Rightarrow |f(x)-\ell|<\varepsilon$ ; la limite est unique.
  \item Outils de calcul : opérations, composition, gendarmes, $|f - \ell| \leq g \to 0$ ; partie entière : $x - 1 < \E(x) \leq x$.
  \item Trigonométrie en $0$ : $\frac{\sin x}{x} \to 1$, $\frac{\tan x}{x} \to 1$, $\frac{1 - \cos x}{x^2} \to \frac{1}{2}$ (démonstration par encadrement).
  \item Prolongement par continuité : possible en $a$ dès que la limite en $a$ existe et est finie.
  \item Suites : $u_n \to a$ et $f$ continue en $a$ $\Rightarrow f(u_n) \to f(a)$ ; suite récurrente convergente $\Rightarrow$ la limite est un point fixe ($f(\ell) = \ell$), \emph{après} avoir prouvé la convergence.
  \item TVI : continuité + changement de signe $\Rightarrow$ existence ; + stricte monotonie $\Rightarrow$ unicité ; dichotomie : amplitude $\frac{b-a}{2^n}$ après $n$ étapes.
  \item Bijection : $f$ continue strictement monotone sur $I$ $\Rightarrow$ $f^{-1}$ continue, même monotonie, courbes symétriques par rapport à $y = x$.
  \item Racine $n$-ième : $y = \sqrt[n]{x} \iff y^n = x$ ($x, y \geq 0$) ; $x^n = a$ : $1$ solution si $n$ impair, $0$ ou $2$ si $n$ pair ; $a^{p/q} = \sqrt[q]{a^p}$.
\end{itemize}
\end{tcolorbox}

\end{document}
